\section{Faisceaux de morphismes}\etiquette[10]{IV.sec10}\pageoriginale

\begin{enonce*}{Proposition 10.1}\etiquette[10.1]{IV.prop10.1}
Soient $E$ un topos, $X$ et $Y$ deux objets de $E$; le foncteur
$Z\mapsto \Hom_E(Z\times X,Y)\simeq \Hom_{E_{/Z}}(X_Z, Y_Z)$ est
représentable.
\end{enonce*}

En effet, les limites inductives sont universelles dans $E$ (II
4.3). Le foncteur $Z\mapsto Z\times X$ commute donc aux limites
inductives. Par suite, le foncteur $Z\mapsto \Hom_E(Z\times X,Y)$
transforme les limites inductives de l'argument $Z$ en limites
projectives. Il est donc représentable \sisi{(IV \hyperref[IV.prop1.4]{1.4} et
\hyperref[IV.theo1.2]{1.2})}{(\hyperref[IV.cor1.2.1]{1.2.1})}.

\setcounter{subsection}{1}
\subsection{}\etiquette[10.2]{IV.subsec10.2}
L'objet représentant le foncteur $Z\mapsto \Hom_{E}(Z\times X, Y)$
est noté $\scrHom_{E}(X,Y)$ (ou plus simplement $\scrHom(X,Y)$) et
est appelé le \emph{faisceau des morphismes de $X$ dans $Y$}. C'est
un bifoncteur en $X$ et $Y$. On a donc un isomorphisme
trifonctoriel.
\begin{equation*}
\Hom_E(Z,\scrHom_E(X,Y))\simeq \scrHom_E(Z\times X,
Y){\quoi}.\tag{10.2.1}
\end{equation*}\etiquette[10.2.1]{IV.eq10.2.1}

Il résulte alors de la formule \eqhyperref[IV.eq10.2.1]{10.2.1} que le bifoncteur
$(X,Y)\mapsto \scrHom(X,Y)$ transforme les limites inductives de
l'argument $X$ (\resp les limites projectives de l'argument $Y$) en
limites projectives dans $E$.

\begin{enonce*}{Proposition 10.3}\etiquette[10.3]{IV.prop10.3}
Soit $v:E\rightarrow E'$ un morphisme de topos. Pour un objet
variable $X$ de $E$ et un objet variable $Y$ de $E'$, on a un
isomorphisme bifonctoriel.
\begin{equation*}
v_{\ast}\scrHom_E(v^{\ast}Y,X)\simeq
\scrHom_{E'}(Y,v_{\ast}X){\quoi}.\tag{10.3.1}
\end{equation*}\etiquette[10.3.1]{IV.eq10.3.1}
\end{enonce*}

\noindent \emph{Démonstration}. Pour tout objet $Z$ de $E'$, on a
une suite d'isomorphismes, fonctoriels en tout les arguments:
\begin{alignat*}{4}
&\Hom_{E'}(Z,v_{\ast}\scrHom_E(v^{\ast}Y,X)) &&\simeq
&& \Hom_E(v^{\ast}Z,\scrHom_E(v^{\ast}Y,X)) &&\quad \text{(adjonction)}\\
& \Hom_E(v^{\ast}Z,\scrHom_E(v^{\ast}Y,X)) && \simeq
&& \Hom_E(v^{\ast}Z\times v^{\ast}Y,X) &&\quad \eqhyperref[IV.eq10.2.1]{10.2.1}\\
&\Hom_E(v^{\ast}Z\times v^{\ast}Y,X) &&\simeq &&
\Hom_E(v^{\ast}(Z\times Y),X) &&\quad \text{($v^{\ast}$ exact à
gauche)}\\
&\Hom_E(v^{\ast}(Z\times Y),X) &&\simeq && \Hom_{E'}(Z\times Y,
v_{\ast}X) &&\quad \text{(adjonction)}\\
& \Hom_{E'}(Z\times Y, v_{\ast}X) &&\simeq &&
\Hom_{E'}(Z,\scrHom_{E'}(Y,v_{\ast}X)) &&\quad \eqhyperref[IV.eq10.2.1]{10.2.1}
\end{alignat*}
Les deux membres de \eqhyperref[IV.eq10.3.1]{10.3.1} représentent donc des foncteurs
isomorphes, \cqfd

\begin{enonce*}{Corollaire 10.4}\etiquette[10.4]{IV.cor10.4}
Soient\pageoriginale $C$ un $𝒰$-site, $X$ un
$𝒰$-préfaisceau sur $C$, $Y$ un $𝒰$-faisceau sur
$C$, $𝒱$ un univers contenant $𝒰$ tel que $C$ soit
$𝒱$-petit, $C^{\widehat{\phantom{,}}}_{𝒱}$ le topos
des $𝒱$-préfaisceaux sur $C$, $C^{\sim}_{𝒰}$ le
topos des $𝒰$-faisceaux sur $C$. Le préfaisceau
$\scrHom_{C^{\widehat{\phantom{,}}}_{𝒱}}(X,Y)$ est un
$𝒰$-faisceau. Soit $X\rightarrow \underline{a}X$ le morphisme
canonique de $X$ dans son faisceau associé. Le morphisme
correspondant
$\scrHom_{C^{\widehat{\phantom{,}}}_{𝒱}}(\underline{a}X,Y)\rightarrow
\scrHom_{C^{\widehat{\phantom{,}}}_{𝒱}}(X,Y)$ est un
isomorphisme. On a un isomorphisme canonique
$\scrHom_{C^{\widehat{\phantom{,}}}_{𝒱}}(\underline{a}X,Y)\simeq
\scrHom_{C^{\sim}_{𝒰}}(\underline{a}X,Y)$.
\end{enonce*}

\setcounter{subsection}{4} \setcounter{subsubsection}{0}
\subsubsection{}\etiquette[10.4.1]{IV.subsubsec10.4.1}
Supposons d'abord que $C$ soit un petit site et que
$𝒱=𝒰$. On a alors un morphisme de topos:
$C^{\sim}_{𝒰}\rightarrow
C^{\widehat{\phantom{,}}}_{𝒰}$ (\sisi{IV }{}\hyperref[IV.subsec4.8]{4.8}), d'où
\eqhyperref[IV.eq10.3.1]{10.3.1} les assertions dans ce cas. Pour passer de là au cas
général, on remarque tout d'abord que le foncteur « faisceau
associé »  ne dépend pas de l'univers (\hyperrefext[II][II.prop3.6]{3.6}) et que le
topos des $𝒱$-faisceaux sur $C$ est équivalent au topos des
$𝒱$-faisceaux sur $C^{\sim}_{𝒰}$ (\hyperrefext[III][III.sec4]{4} et
\hyperref[IV.sec1]{1}). Il suffit donc de démontrer le lemme suivant:

\begin{enonce*}{Lemme 10.4.2}\etiquette[10.4.2]{IV.lem10.4.2}
Soient $E$ un $𝒰$-topos, $𝒱$ un univers contenant
$𝒰$, $E^{\sim}_{𝒱}$ le topos des
$𝒱$-faisceaux sur $E$, $X$ et $Y$ deux objets de $E$. Il
existe un isomorphisme canonique $\scrHom_E(X,Y)\simeq
\scrHom_{E^{\sim}_{V}}(X,Y)$.
\end{enonce*}

\setcounter{subsubsection}{2}
\subsubsection{}\etiquette[10.4.3]{IV.subsubsec10.4.3}
Le foncteur $\epsilon_{E}:E\rightarrow E_{\widetilde{𝒱}}$
est pleinement fidèle et exact à gauche. Par suite on a, pour tout
objet $Z$ de $E$ un isomorphisme
$$
\Hom_{E_{\widetilde{𝒱}}}(\epsilon_E Z, \epsilon_E
\scrHom_E(X,Y))\simeq \Hom_{E_{\widetilde{𝒱}}}(\epsilon_E
Z\times \epsilon_E X, \epsilon_E Y)
$$
Mais tout objet de $E_{\widetilde{𝒱}}$ est limite inductive
d'objets provenant de $E$, d'où \eqhyperref[IV.eq10.2.1]{10.2.1} l'assertion.

\subsection{}\etiquette[10.5]{IV.subsec10.5}
Soit $v:E\rightarrow E'$ un morphisme de topos. On se propose de
définir quatre morphismes bifonctoriels.
\begin{alignat*}{4}
\Phi_v & : v^{\ast}\scrHom(X,Y) &&\longrightarrow
\scrHom(v^{\ast}X,v^{\ast}Y) &\quav
& X \text{ et } Y \text{ objets de } E'{\quoi},\\
\Psi_v & : v_{\ast}\scrHom(X,Y) &&\longrightarrow \scrHom(v_{\ast}X,
v_{\ast}Y) &\quav & X \text{ et } Y \text{ objets de } E{\quoi},\\
\Xi_v & : \scrHom(v_{!}X,Y) &&\longrightarrow
v_{\ast}\scrHom(X,v^{\ast}Y) &\quav & X\in \ob E, Y\in \ob
E'{\quoi},\\
\Lambda_v & : v_{!}(X\times v^{\ast}Y) &&\longrightarrow
v_{!}(X)\times Y & & X\in \ob E, Y\in \ob E'{\quoi},
\end{alignat*}
où,\pageoriginale pour définir les morphismes $\Xi_v$ et $\Lambda_v$
on suppose que le foncteur $v^{\ast}$ image réciproque par $v$ admet
un adjoint à gauche $v_{!}$.

\subsubsection{}\etiquette[10.5.1]{IV.subsubsec10.5.1}
\emph{Définition de} $\Psi_v$. Soit $X$ un objet de $E$. On a un
morphisme d'adjonction $v^{\ast}v_{\ast}X\rightarrow X$, d'où, pour
tout objet $Z$ de $E'$, un morphisme bifonctoriel $v^{\ast}(Z\times
v_{\ast}X)\rightarrow (v^{\ast}Z)\times X$. On a donc, pour tout
objet $Y$ de $E$, un morphisme trifonctoriel
$$
\Hom_E((v^{\ast}Z)\times X, Y)\longrightarrow
\Hom_E(v^{\ast}(Z\times v_{\ast}X),Y){\quoi}.
$$
On en déduit, par adjonction, un morphisme trifonctoriel
$$
\Hom_{E'}(Z,v_{\ast}\scrHom(X,Y))\longrightarrow \Hom_{E'}(Z,
\scrHom(v_{\ast}X, v_{\ast}Y)){\quoi};
$$
d'où le morphisme $\Psi_v$.

\subsubsection{}\etiquette[10.5.2]{IV.subsubsec10.5.2}
\emph{Définition de} $\Lambda_v$. Pour tout objet $X$ de $E$, on a
un morphisme d'adjonction $X\rightarrow v^{\ast}v_{!}X$, d'où, pour
tout objet $Y$ de $E$, un morphisme bifonctoriel $X\times
v^{\ast}Y\rightarrow v^{\ast}v_{!}(X)\times v^{\ast}Y\simeq
v^{\ast}(v_{!}(X)\times Y)$; d'où, par adjonction, le morphisme
$\Lambda_v$.

\subsubsection{}\etiquette[10.5.3]{IV.subsubsec10.5.3}
\emph{Définition de} $\Xi_v$. Soient $X$ un objet de $E$, $Z$ et $Y$
deux objets de $E'$. Le morphisme $\Lambda_v:v_{!}(X\times
v^{\ast}Z)\rightarrow v_{!}(X)\times Z$ fournit un morphisme
trifonctoriel $\Hom_{E'}(v_{!}(X)\times Z,Y)\rightarrow
\Hom_{E'}(v_{!}(X\times v^{\ast}Z),Y))$; d'où, par adjonction, un
morphisme trifonctoriel
$$
\Hom_{E'}(Z,\scrHom(v_{!}X,Y))\longrightarrow
\Hom_{E'}(Z,v_{\ast}\scrHom(X,v^{\ast}Y)){\quoi};
$$
d'où le morphisme $\Xi_v$.

\subsubsection{}\etiquette[10.5.4]{IV.subsubsec10.5.4}
\emph{Définition de} $\Phi_v$. Le morphisme trifonctoriel de
(\hyperref[IV.subsubsec10.5.3]{10.5.3})
$$
\Hom_{E'}(v_{!}(Z)\times X,Y)\longrightarrow \Hom_{E'}(v_{!}(Z\times
v^{\ast}X),Y){\quoi};
$$
où $Z$ est un objet de $E$ et $X$, $Y$ sont des objets de $E'$,
permet d'obtenir, par adjonction, un morphisme trifonctoriel,
$$
\Hom_E(Z,v^{\ast}\scrHom(X,Y))\longrightarrow
\Hom_E(Z,\scrHom(v^{\ast}X, v^{\ast}Y)){\quoi};
$$
d'où une définition du morphisme $\Phi_v$. Il existe une deuxième
manière de définir $\Phi_v$. Soient $X$, $Y$, $Z$ trois objets de
$E'$. Le foncteur $v^{\ast}:E'\rightarrow E$ fournit un morphisme
trifonctoriel
$$
\Hom_{E'}(Z\times X,Y)\longrightarrow \Hom_E (v^{\ast}(Z)\times
v^{\ast}(X), v^{\ast}Y){\quoi};
$$
d'où, par adjonction, un morphisme trifonctoriel
$$
\Hom_{E'}(Z,\scrHom(X,Y))\longrightarrow
\Hom_{E'}(Z,v_{\ast}\scrHom(v^{\ast}X, v^{\ast}Y)){\quoi}.
$$
On\pageoriginale a donc un morphisme bifonctoriel
$$
\scrHom(X,Y)\longrightarrow v_{\ast}\scrHom(v^{\ast}X,
v^{\ast}Y){\quoi};
$$
d'où, par adjonction, un morphisme dont on laisse au lecteur le soin
de vérifier que c'est bien le morphisme $\Phi_v$ défini ci-dessus.

\begin{enonce*}{Proposition 10.6}\etiquette[10.6]{IV.prop10.6}
Soit $v:E\rightarrow E'$ un morphisme de topos.
\begin{enumerate}
\renewcommand{\labelenumi}{\theenumi)}
\item Le morphisme $\Psi_v$ $(\hyperref[IV.subsubsec10.5.1]{10.5.1})$ est un
  isomorphisme pour tout
objet $Y$ de $E$ si et seulement si le morphisme d'adjonction
$v^{\ast}v_{\ast}X\rightarrow X$ est un isomorphisme. En
particulier, le morphisme $\Psi_v$ est un isomorphisme pour tout
couple d'objets $(X,Y)$ si et seulement si le foncteur
$v_{\ast}:E\rightarrow E'$ est pleinement fidèle, \ie si $v$ est un
plongement de topos.

\item Les conditions suivantes sont équivalentes:
\begin{enumerate}
\renewcommand{\theenumii}{\roman{enumii}}
\renewcommand{\labelenumii}{(\theenumii)}
\item Pour tout couple $(X,Y)$ d'objets de $E'$, le morphisme $\Phi_v$
$(\hyperref[IV.subsubsec10.5.4]{10.5.4})$ est un isomorphisme.

\item Le foncteur $v^{\ast}$ admet un adjoint à gauche $v_{!}$ et pour
tout objet $X$ de $E$ et tout objet $Y$ de $E'$, le morphisme
$\Xi_v$ $(\hyperref[IV.subsubsec10.5.3]{10.5.3})$ est un isomorphisme.

\item Le foncteur $v^{\ast}$ admet un adjoint à gauche $v_{!}$ et pour
tout objet $X$ de $E$ et pour tout objet $Y$ de $E'$, le morphisme
$\Lambda_v$ $(\hyperref[IV.subsubsec10.5.2]{10.5.2})$ est un isomorphisme.
\end{enumerate}
\end{enumerate}
\end{enonce*}

\setcounter{subsection}{6} \setcounter{subsubsection}{0}
\subsubsection{}\etiquette[10.6.1]{IV.subsubsec10.6.1}
La première assertion résulte immédiatement de
(\hyperref[IV.subsubsec10.5.1]{10.5.1}). Il résulte aussi immédiatement de
(\hyperref[IV.subsubsec10.5.2]{10.5.2}), (\hyperref[IV.subsubsec10.5.3]{10.5.3}),
(\hyperref[IV.subsubsec10.5.4]{10.5.4}) que (iii) $\Leftrightarrow$ (ii)
$\Leftrightarrow$ ((i) et le foncteur $v^{\ast}$ admet un adjoint à
gauche $v_{!}$). Il reste donc à montrer que (i) implique que
$v^{\ast}$ admet un adjoint à gauche, \ie (\hyperref[IV.rem1.8]{1.8}) que (i)
implique que $v^{\ast}$ commute aux petits produits. Soient $e'$
l'objet final de $E'$ et $I$ un petit ensemble. Comme $v^{\ast}$ est
exact à gauche, $v^{\ast}(e')$ est un objet final de $E$ et comme
$v^{\ast}$ commute aux limites inductives $v^{\ast}(\amalg_I
e')\simeq \amalg_I v^{\ast}(e')$. Pour tout objet $Y$ de $E'$, on a
$\scrHom(\amalg_I e', Y)\simeq \Pi_I\scrHom(e',Y)\simeq \Pi_I Y$ et
$\scrHom(\amalg_I v^{\ast}(e')$, $v^{\ast}Y)\simeq
\Pi_i\scrHom(v^{\ast}(e'), v^{\ast}Y)\simeq \Pi_I
v^{\ast}Y$.\pageoriginale Le morphisme fonctoriel $\Phi_v$ induit
donc un isomorphisme $v^{\ast}\Pi_I Y\simeq \Pi_I v^{\ast}Y$ dont on
vérifie que c'est l'isomorphisme canonique. \cqfd

\begin{enonce*}{Corollaire 10.7}\etiquette[10.7]{IV.cor10.7}
Soient $E$ un topos, $Z$ un objet de $E$, $j_Z:E_{/Z}\rightarrow E$
le morphisme de localisation (\hyperref[IV.subsec5.2]{5.2}), $X$, $Y$
deux objets de $E$.
\begin{enumerate}
\renewcommand{\labelenumi}{\theenumi)}
\item Il existe un isomorphisme bifonctoriel
$$
j^{\ast}_Z\scrHom(X,Y)\simeq \scrHom(j^{\ast}_ZX,
j^{\ast}_ZY){\quoi}.
$$

\item Soit $X'$ un objet de $E_{/Z}$. Il existe un isomorphisme
bifonctoriel
$$
\scrHom_E(j_{Z!}X',Y)\simeq j_{Z\ast}\scrHom_{E/Z}(X',j^{\ast}_Z
Y){\quoi}.
$$

\item Soit $Y'$ un objet de $E/Z$. Lorsque $Z$ est un ouvert de $E$,
il existe un isomorphisme bifonctoriel
$$
j_{Z\ast}\scrHom(X',Y')\simeq \scrHom(j_{Z\ast}X',
j_{Z\ast}Y'){\quoi}.
$$
\end{enumerate}
\end{enonce*}

Les assertions 1) et 2) se démontrent en remarquant que le morphisme
$\Lambda_{j_Z}$ est un isomorphisme. Pour l'assertion 3), il suffit
de remarquer que $j_{Z\ast}$ est pleinement fidèle, car $j_{Z!}$ est
pleinement fidèle (\hyperrefext[I][I.rem5.7]{5.7}. a)).

\begin{enonce*}{Corollaire 10.8}\etiquette[10.8]{IV.cor10.8}
Soient $E$ un topos, $X$, $Y$ et $Z$ trois objets de $E$,
$j_Z:E_{/Z}\rightarrow E$ le morphisme de localisation.
\begin{enumerate}
\renewcommand{\labelenumi}{\theenumi)}
\item On a un isomorphisme canonique
$$
j_{Z\ast}j^{\ast}_ZY\simeq \scrHom(Z,Y){\quoi}.
$$

\item On a un isomorphisme canonique
$$
\Hom_E(Z,\scrHom(X,Y))\simeq \Hom_{E_{/Z}}(j^{\ast}_ZX,
j^{\ast}_ZY){\quoi}.
$$
\end{enumerate}
\end{enonce*}

Soit $e_Z$ l'objet final de $E_{/Z}$ (\ie l'objet
$\id_Z:Z\rightarrow Z$)). Il résulte de \eqhyperref[IV.eq10.2.1]{10.2.1} qu'on a un
isomorphisme canonique
$$
\scrHom_{E/Z}(e_Z, j^{\ast}_ZY)\simeq j^{\ast}_ZY{\quoi};
$$
d'où, d'après (\hyperref[IV.cor10.7]{10.7} 2), un isomorphisme
$$
j_{Z\ast}j^{\ast}_Z Y\simeq \scrHom(j_{Z!}e_Z, Y){\quoi.}
$$
Comme $j_{Z!}e_Z=Z$, on a démontré 1). Démontrons 2). De
(\hyperref[IV.cor10.7]{10.7}), on tire un\pageoriginale isomorphisme canonique
$$
\Hom_{E_{/Z}}(e_Z, j^{\ast}_Z\scrHom(X,Y))\simeq
\Hom_{E_{/Z}}(j^{\ast}_Z X, j^{\ast}_ZY){\quoi};
$$
d'où, par adjonction sur le premier membre,
$$
\Hom_E(Z,\scrHom(X,Y))\simeq \Hom_{E_{/Z}}(j^{\ast}_ZX,
j^{\ast}_ZY){\quoi}.
$$
